\chapter[Instâncias de Teste]{Instâncias de Teste}
\label{ape:instancias}

Este apêndice descreve as 30 instâncias sintéticas utilizadas nos experimentos, incluindo parâmetros de geração, estrutura dos arquivos e um exemplo representativo.

\section{Visão Geral}

As instâncias foram geradas sinteticamente para o TSP-SD-ATP, variando o número de pontos de interesse ($n$) e a disposição espacial (variantes a, b, c). A \autoref{tab:ape-instancias} lista todas as instâncias e suas características.

\begin{table}[htbp]
\centering
\caption{Instâncias de teste.}
\label{tab:ape-instancias}
\begin{tabular}{clrr}
\toprule
Grupo & Instâncias & $n$ & Variantes \\
\midrule
Pequenas & 10a, 10b, 10c & 10 & a, b, c \\
Pequenas & 11a, 11b, 11c & 11 & a, b, c \\
Pequenas & 12a, 12b, 12c & 12 & a, b, c \\
Pequenas & 13a, 13b, 13c & 13 & a, b, c \\
Pequenas & 14a, 14b, 14c & 14 & a, b, c \\
Pequenas & 15a, 15b, 15c & 15 & a, b, c \\
Médias & 30a, 30b, 30c & 30 & a, b, c \\
Grandes & 50a, 50b, 50c & 50 & a, b, c \\
Grandes & 100a, 100b, 100c & 100 & a, b, c \\
\bottomrule
\end{tabular}
\legend{Fonte: dados da pesquisa.}
\end{table}

\section{Parâmetros de Geração}

Cada instância foi gerada com uma semente pseudoaleatória, resultando em coordenadas $(x, y)$ distribuídas uniformemente no plano. O tensor de custo tridimensional $G[i][j][k]$ é calculado como:

\begin{equation}
G[i][j][k] = d[j][k] + \alpha \cdot \theta(i,j,k)
\end{equation}

onde $d[j][k]$ é a distância euclidiana entre $j$ e $k$, e $\theta(i,j,k)$ é a penalidade angular dependente da tripla $(i,j,k)$.

\section{Exemplo: Instância 10a}

A \autoref{tab:ape-coords-10a} apresenta as coordenadas dos 10 pontos de interesse da instância 10a, incluindo a base (ponto 0).

\begin{table}[htbp]
\centering
\caption{Coordenadas da instância 10a.}
\label{tab:ape-coords-10a}
\begin{tabular}{crr}
\toprule
Ponto & $x$ & $y$ \\
\midrule
0 (base) & 50,00 & 50,00 \\
1 & 22,21 & 33,09 \\
2 & 87,13 & 49,90 \\
3 & 12,79 & 85,09 \\
4 & 63,99 & 14,89 \\
5 & 78,41 & 67,97 \\
6 & 32,87 & 42,95 \\
7 & 55,42 & 91,77 \\
8 & 43,75 & 55,79 \\
9 & 91,25 & 27,63 \\
\bottomrule
\end{tabular}
\legend{Fonte: dados da pesquisa.}
\end{table}
